\section{第\ref{sec:motorlearning}章　運動学習}

最小二乗則と専用の誤差配線の利点は定理である。誤差の配線を持つ回路の構成、一致による弱化、遅延へのトレースの調整、後効果と技能の自発的回復は測定されている。回路全体が教師あり学習として働き、内部モデルを作るというのは有力な仮説である。二過程モデルの数値は本書のモデルによる計算である。

{\footnotesize
\setlength{\tabcolsep}{3pt}\renewcommand{\arraystretch}{1.15}
\begin{longtable}{>{\raggedright}p{0.5cm}>{\raggedright}p{4.0cm}>{\raggedright}p{5.6cm}>{\raggedright}p{2.3cm}>{\raggedright\arraybackslash}p{1.7cm}}
\toprule
No. & 主張 & 実験と測定されたもの & 出典 & 状態 \\
\midrule\endfirsthead
\toprule
No. & 主張 & 実験と測定されたもの & 出典 & 状態 \\
\midrule\endhead
1 & 規則~\eqref{eq:lms}は平均すると二乗誤差の勾配をたどる & 適応フィルタ理論の結果：入力と出力の誤差に比例する補正は、平均二乗誤差の確率的勾配降下である。 & \cite{WidrowHoff1960} & 理論 \\
2 & 出力ごとに誤差の配線があれば速さは出力の数によらない；一つの共通の数では、ノイズを出すニューロンの数とともに落ちる（命題~\ref{prop:teachers}） & 線形ネットワークの学習曲線：ニューロンのランダムな摂動による学習は、摂動を受けるニューロンの数の倍だけ正確な勾配より遅い。 & \cite{Werfel2005} & 理論 \\
3 & 誤差の配線を持つ回路は教師あり学習として働く（命題~\ref{prop:errorwire}） & 回路の構造から導かれた理論。7--10行の実験はそれと整合する；回路全体がまさにそう働くことは証明されていない。 & \cite{Marr1969, Albus1971} & 仮説 \\
4 & 拡張層：約 $7\cdot10^{10}$ 個の小さな\nt{GLUT}ニューロン、脳の残り全部より多い & 溶解した組織の懸濁液での細胞核の計数：ヒトの小脳には脳全体の $86\cdot10^9$ のうち $69\cdot10^9$ のニューロン。立体解析学：高齢男性の小脳に $101\cdot10^9$ 個の小さなニューロン。 & \cite{Azevedo2009, Andersen1992cereb} & 測定済み \\
5 & 入力の次元を上げると、望みの関係が線形になる & 分割の数についての定理：$n$ 個の入力の重み付き和としきい値は、一般の位置にある約 $2n$ 個のベクトルの任意のラベル付けを実現する。小脳がまさにこれを使っているというのは3行目の仮説の一部。 & \cite{Cover1965} & 理論 \\
6 & 約 $3\cdot10^7$ 個の出力\nt{GABA}ニューロン、それぞれ $10^5$ 程度の入力 & ヒト小脳の立体解析学：出力ニューロンは $30.5\cdot10^6$ 個。電子顕微鏡、ラット：出力ニューロン1個あたり拡張層からの入力は約 $1.75\cdot10^5$。出力ニューロンの抑制性伝達物質は総説による。 & \cite{Andersen1992cereb, NapperHarvey1988, Ito2001} & 測定済み \\
7 & 出力ニューロンあたり誤差の配線はちょうど1本、毎秒1回ほど、毎回強力なバースト & 記録の総説：各出力ニューロンは1本の登上\nt{GLUT}入力を受け、それが $1$~Hz 程度の頻度で複雑スパイクを起こす。 & \cite{Ito2001} & 測定済み \\
8 & バーストの確率は運動誤差の方向に依存する & サル、小脳の眼球運動領域、サッカード適応：視覚誤差の方向が複雑スパイクの確率を、大きさがそのタイミングを決める；バーストは次の試行の単純スパイクを変え、細胞の好みの誤差の方向に運動をずらす。 & \cite{Herzfeld2018} & 測定済み \\
9 & 誤差のバーストと活動中の入力が一致すると、その入力が弱まる（$-\eta e_i x_j$） & 実験の総説：拡張層からの入力と誤差の配線を同時に刺激すると、その入力が長期的に弱まる；入力だけの刺激では弱まらない。 & \cite{Ito2001} & 測定済み \\
10 & トレースの長さは誤差の遅延に合わせてある：遅延 $100\ms$ では $100\ms$ 前に活動していた入力が最も強く弱まる & スライスと生きたマウス：眼球運動学習を担う領域では、弱化は入力と誤差の間隔約 $120\ms$（この課題で誤差信号が届くのにかかる時間）に鋭く同調している；別の領域では細胞ごとに異なる間隔に同調している。 & \cite{Suvrathan2016} & 測定済み \\
11 & 回路は体の内部モデルを学習する適応フィルタ~\eqref{eq:adaptivefilter}である & 回路の構造から導かれたモデル。学習したフィルタを体の伝達関数として測定した直接の実験はない。 & \cite{Fujita1982} & 仮説 \\
12 & 予測は自分の動きの結果を差し引くので、自分をくすぐることはできない & ヒト：自分で触れると、他人の同じ接触よりくすぐったくない；fMRIでは体性感覚皮質の自分の接触への応答が弱く、動きが皮膚に触れるとき小脳の活動が少ない。予測の差し引きによる説明は仮説。 & \cite{Blakemore1998} & 仮説 \\
13 & 外力のもとで外れは小さくなり、外力を取り除くと手は反対側に外れる & ヒトが力場の中でロボットのハンドルを持って手を伸ばす：軌道はまっすぐに戻る；力場を取り除いた後の後効果は最初のゆがみのほぼ鏡像で、訓練していない空間の部分にも転移する。 & \cite{ShadmehrMussaIvaldi1994} & 測定済み \\
14 & 二つの過程が学習する~\eqref{eq:tworate}；逆向きの外乱の後、技能はひとりでに戻る & ヒト、力場、次に短い逆向きの力場と誤差を固定した試行：補正はひとりでに古い値へ戻る；二過程モデルはこれを再現し、一過程モデルは再現しない。 & \cite{Smith2006} & 測定済み \\
15 & 図~\ref{fig:tworate}の数値：$200$ 回の動きで $0.84$（遅い過程 $0.63$）；逆向き $6$ 回；速い $-0.53$、遅い $0.46$；$40$ 回で $0.37$ まで回復 & \texttt{models/\allowbreak motor\_\allowbreak learning.py} による計算、$A_f=0.92$, $B_f=0.1$, $A_s=0.996$, $B_s=0.02$：$0.840$（$0.634$ と $0.205$）、$6$ 回、$-0.531$ と $0.457$、$40$ 回目でピーク $0.371$。 & \texttt{models/\allowbreak motor\_\allowbreak learning.py} & モデル \\
16 & 体が異なる速さで変わるので、二つの過程が必要になる & 理論：寿命の異なる複数の外乱のもとでの最適推定は、時定数の異なる複数のフィルタを与え、自発的回復と再学習の加速を説明する。 & \cite{Kording2007} & 仮説 \\
\bottomrule
\end{longtable}}
